\section{Discussion and limitations}\label{sec:limits}
\paragraph{Scope of the model}
All results are from simulation of a reduced model. The model is one leg with two drives, its share of the trunk on a vertical rail, and a quarter share of the battery pack, which assumes that the four legs of a pronking quadruped load the pack equally. The drives are averaged models with $i_d=0$ and an ideal current loop; the validation against a switching-level model (Section~\ref{sec:model}) quantifies the error of this choice for one drive, and the plant-only current-lag variant tests the filter against a slower current loop (Table~\ref{tab:abl}). The brake chopper, not the filter, handles overvoltage. Several parameters are assumptions rather than measured or sourced values, among them the phase resistance, the drives' idle power, the chopper, the auxiliary load and the sensor noise; Table~\ref{tab:params} marks each of them.

\paragraph{Battery parameters in the filter}
The guarantee depends on knowing the battery's electromotive force $E$ and series resistance $R_0$ exactly. The power budget $\bar P$ is computed from both, and the one-step guarantee (Corollary~\ref{cor:closed}) holds only if they are exact. The true-state filter takes $E$ from the simulated battery and $R_0$ from the model; the EKF filter estimates both from the measured bus voltage and battery current (Section~\ref{sec:estimator}). The bus capacitance, the substep, the chopper, the drive power model and the sensors' noise levels are still assumed known. The guarantee does not extend to estimated values, and the margins, fixed or confidence-scaled, are not derived from a bound on the estimation error. The first-order rule of Remark~\ref{rem:sens} bounds every observed overshoot of the true-state filter under estimation noise only for error bounds of at least \res{sens_kall} times the noise scale, and with one noise scale it exceeds the fixed current margin $\delta_I=\res{env_margini}$~A in all \res{sens_ncases} cases (Table~\ref{tab:sens}).

Estimation helps against a wrong resistance but not against unmodelled sensor noise (Section~\ref{sec:estresults}). The confidence took only values near one or zero, so the margin acted as a switch, and we have not shown a benefit of scaling it continuously. The policies observe the plant's bus voltage and state of charge in every evaluation, also when the filter uses estimates. Finally, the filter uses the same backward-Euler bus equation as the simulated plant. Zero violations with exact parameters therefore show that the filter is consistent with its own model, not that it is robust to other model errors.

\paragraph{Training}
Two training issues qualify the comparison. First, trained from scratch with the filter for \res{train_steps_scratch} steps, only one of \res{n_seeds} seeds learned to hop repeatedly (Table~\ref{S-tab:perseed}); the hopping filter-trained policies of the main comparison fine-tune the pretrained baseline policies with the filter in the loop. Every method was trained with \res{n_seeds} seeds, except the penalty weights \res{pen_ws_w_part}, trained with one seed each. This limits the statistical strength of the comparisons, and several findings differ between seeds (Tables~\ref{S-tab:perseed} and~\ref{S-tab:perseedws} of the supplementary material). Second, the SAC-Lagrangian baseline did not meet its cost limit within its \res{train_steps_ws} training steps, and its multiplier was still at its largest value at the end of training (Table~\ref{tab:lag}). The comparison with constrained reinforcement learning is therefore with this training budget and the traditional multiplier update, not with a converged constrained policy. Proportional-integral-derivative (PID) Lagrangian methods, which add proportional and derivative terms to the multiplier update \cite{stooke2020pid}, are a stronger constrained baseline that we did not test.

\paragraph{Performance cost of the limits}
Meeting the limits costs height, and the filter-trained policies do not reach the largest height within the limits. The largest heights in the main evaluation come with breaches and so do not show what can be reached within the limits (Table~\ref{S-tab:wsfull} of the supplementary material and Figure~\ref{fig:tradeoff}). Within the limits and without falls, the warm-started penalty policies with the test-time filter reached a larger mean max height than the filter-trained policies at every state of charge, for each weight trained with \res{n_seeds} seeds ($w$ of \res{pen_ws_w_full}; Table~\ref{S-tab:wsfull}). In these evaluations, then, training with the filter cost height against a penalty policy filtered at test time, and its gain was the absence of the stress-test falls (Section~\ref{sec:testtime}), which came from one of \res{n_seeds} penalty seeds.

\paragraph{Battery charge limit}
The budget does not model the charge limit $I_{\mathrm{chg}}$. When the barrier conditions can only be met by regenerating more current than the battery accepts, the plant clamps the charge current and the bus voltage ends higher than the filter predicts. The filter is then conservative rather than unsafe, but its budget no longer describes the admissible loads exactly, and the brake chopper, not the filter, handles the resulting rise in bus voltage.

\paragraph{Unrealisable drive currents in the baselines}
Without the filter, the drives are occasionally asked for currents the inverter cannot realise. Assumption (A6) fails when a drive's voltage circle and its current limits do not intersect, and the drive model then applies the nearest current in the current box. In the main evaluations this happened in \res{asix_base_sub} substeps over \res{asix_base_ep} episodes in baseline mode (\res{asix_base_ppm} per million substeps) and in \res{asix_pen_sub} substeps over \res{asix_pen_ep} episodes in penalty mode (\res{asix_pen_ppm} per million). It never happened in filter mode (\res{asix_filt_ep} episodes) or in the Lagrangian policies' own mode (\res{asix_lag_ep} episodes). The baseline policies under the parameter variants add \res{asix_var_sub} substeps. It occurs in the substep after a bus collapse: we checked \res{asixchk_ep} episodes of the penalty policy with the most occurrences, one per state of charge, and of its \res{asixchk_sub} such substeps, \res{asixchk_after} came directly after a substep in which the bus collapsed and \res{asixchk_uv} started with the bus voltage below $V_{\min}$.

\paragraph{Computation}
Real-time execution on robot hardware is not shown. The filter is implemented in Python; one call took \res{filt_us_mean}~\textmu s on average and \res{filt_us_p99}~\textmu s at the 99th percentile on the training CPU (Section~\ref{sec:filter}), against a substep of \res{env_dtsim_ms}~ms. The work per call is bounded in advance: the budget is evaluated in closed form and the projection takes at most \res{filt_maxdoublings} doublings and \res{filt_halvings} halvings of the multiplier bracket (Algorithm~\ref{alg:filter}).

\paragraph{Future work}
The next steps are experiments on hardware with measured bus voltage and battery current; a full quadruped in which all legs share one bus, with gaits other than pronking; thermal limits of the drives alongside the electrical ones \cite{wan2026thermal,qian2026thermal}; estimators with noise models matched to the sensors and margins derived from bounds on their errors, for example through the rule of Remark~\ref{rem:sens}; and a compiled implementation of the filter.
